\documentclass[12pt]{article}
\usepackage[utf8]{inputenc}
\usepackage[T1]{fontenc}
\usepackage[margin=1in]{geometry}
\usepackage{setspace}
\usepackage{xcolor}
\usepackage[colorlinks=true, urlcolor=blue!60!black, linkcolor=blue!60!black]{hyperref}
\setlength{\parindent}{2em}
\setlength{\parskip}{1.25em}
\renewcommand{\baselinestretch}{1.0}
\begin{document}
\begin{center}
{\LARGE \textit{Fletcher v. Peck}, 10 U.S. 87 (1810)\par}
\vspace{0.8em}
{\large Case Brief --- Four-Page Edition\par}
\end{center}
\vspace{1em}

\section{Memory Jogger}

In the first Supreme Court decision invalidating a state statute as unconstitutional, Chief Justice Marshall held that a state's legislative land grant is an executed ``contract'' under the Contract Clause, barring Georgia from rescinding its 1795 Yazoo land grant against innocent third-party purchasers even though the original statute was procured by massive legislative bribery (\href{https://www.courtlistener.com/opinion/84935/fletcher-v-peck/}{10 U.S. 87}).

\section{Detailed Case Facts}

In 1795 the Georgia General Assembly sold about 35 million acres along the Yazoo River to four land companies for roughly 1.5 cents per acre. Company representatives had promised legislators shares in the land for their votes. Public outrage swept most of them from office, and in 1796 the new legislature declared the grant null and void and ordered its records expunged. By then the grantees had resold large tracts to northern speculators and innocent purchasers without notice of the bribery.

In 1803 John Peck conveyed 15,000 acres to Robert Fletcher for \$3,000 with covenants that Georgia had authority to sell, that the 1795 act was valid, that the 1796 act had not impaired title, and that Georgia was seised in fee simple despite unextinguished Native American occupancy. Because the Eleventh Amendment barred suing Georgia directly, the parties arranged what Justice Johnson recognized as a ``friendly'' or ``feigned'' test case.

\section{Procedural History}

Fletcher sued Peck for breach of covenant in the federal circuit court in Massachusetts on four counts matching the covenants. After a delay while Congress considered compensating Yazoo claimants, Peck pleaded that he and intervening purchasers bought for value without notice. The circuit court gave judgment for Peck on all four counts. On writ of error, the Supreme Court heard argument in 1809, reargument in 1810, and affirmed on March 16, 1810.

\section{Judicial Votes}

Affirmed, \textbf{4--1} on the reasoning (unanimous on the rescission issue), with five Justices participating.
\begin{itemize}
\item \textbf{Majority (4):} Chief Justice Marshall, joined by Justices Washington, Livingston, and Todd.
\item \textbf{Concurring in part and dissenting in part:} Justice William Johnson, who agreed Georgia could not revoke its grant but relied on natural justice rather than the Contract Clause, and dissented on the Indian-title issue.
\item \textbf{Not participating:} Justices Cushing and Chase, both ill.
\end{itemize}

\section{Holding}

A completed state land grant is an executed contract under Article~I, \S~10. Once title has vested in bona fide purchasers for value without notice, a later legislature cannot rescind the grant, and a court in a collateral private suit will not void a duly enacted statute by probing legislators' corrupt motives. The 1795 act was valid under Georgia's 1789 constitution, and a state may hold a fee-simple estate subject to unextinguished Indian occupancy. Marshall's appeal to ``general principles which are common to our free institutions,'' his bill-of-attainder analogy, and his remark about rescission against the original corrupt grantees are dicta.

\section{Analysis of Opinions}

\textbf{Majority (Marshall).} Marshall moved count by count. Declaring a law void is ``at all times, a question of much delicacy,'' and nothing in Georgia's constitution barred the sale. Courts cannot feasibly probe individual legislators' motives or count tainted votes, especially with Georgia absent from the record, and equity protects bona fide purchasers. On rescission, ``when a law is in its nature a contract, when absolute rights have vested under that contract, a repeal of the law cannot devest those rights.'' A grant is an executed contract that extinguishes the grantor's right and implies an obligation not to reassert it, and the Clause's text draws no line between private and public contracts. The Framers designed Article~I, \S~10 as a ``bill of rights for the people of each state.'' Finally, Indian title was a right of occupancy compatible with Georgia's fee simple.

\textbf{Johnson (concurring in part, dissenting in part).} Johnson rested on natural law, ``a principle which will impose laws even on the deity,'' warning that extending the Contract Clause to executed grants went ``very far beyond the obvious and necessary import of the words'' and could paralyze limitations statutes, recording acts, and eminent domain. He would hold that Georgia had at most a future interest before Indian title was extinguished, and he doubted the Court should hear a case bearing ``strong evidence, upon the face of it, of being a mere feigned case.''

\section{Examples: Future Applications}

\begin{enumerate}
\item \textbf{Same-Side (Unconstitutional): University Charter.} A later legislature converts a privately chartered, tax-exempt university into a state institution without the trustees' consent. A state charter, like a land grant, is an executed contract (mirroring \textit{Dartmouth College}).
\item \textbf{Same-Side (Unconstitutional): Curative Tax-Sale Repeal.} After backlash, a legislature retroactively repeals a curative statute to strip title from bona fide buyers of tax-forfeited land, destroying vested rights (mirroring \textit{Wood v. Lovett}).
\item \textbf{Opposite-Side (Constitutional): Solar Subsidy.} A legislature prospectively repeals an annual subsidy to solar manufacturers that created no binding contract or conveyance. Unvested benefits may be repealed.
\item \textbf{Opposite-Side (Constitutional): Mortgage Moratorium.} A temporary, narrowly tailored Depression-era moratorium requiring reasonable rent to lenders is a valid police-power measure (as in \textit{Blaisdell}).
\item \textbf{Fence-Sitter: Bribed Mineral Lease.} A legislature rescinds a bribe-procured mineral lease and sues in equity before any transfer to innocent parties. Marshall left this open, but \textit{Fletcher}'s reluctance to probe legislative motive creates friction.
\end{enumerate}

\section{Critique}

\textbf{Originalist and textualist views.} The Contract Clause was aimed chiefly at state debtor-relief laws affecting private contracts, and the Committee of Style dropped ``private'' without recorded debate. Echoing Johnson, textualists argue an executed grant is a completed conveyance, not an ongoing ``obligation of contracts,'' and fault Marshall's appeal to unwritten principles and the Court's willingness to hear a collusive suit.

\textbf{Progressive and doctrinal views.} Realist historians see early judicial overreach that shielded a bribed giveaway of 35 million acres from democratic reversal, paving the way for \textit{Dartmouth College} and economic substantive due process. Federal Indian law scholars criticize the cursory treatment of indigenous ownership as a step toward \textit{Johnson v. M'Intosh}.

\section{Key Quotations}

\begin{quote}
``A contract executed, as well as one which is executory, contains obligations binding on the parties. A grant, in its own nature, amounts to an extinguishment of the right of the grantor, and implies a contract not to reassert that right. A party is, therefore, always estopped by his own grant.''
\end{quote}
\noindent --- Chief Justice Marshall, Opinion of the Court, 10 U.S. at 137.

\begin{quote}
``It is, then, the unanimous opinion of the court, that, in this case, the estate having passed into the hands of a purchaser for a valuable consideration, without notice, the state of Georgia was restrained, either by general principles which are common to our free institutions, or by the particular provisions of the constitution of the United States, from passing a law whereby the estate of the plaintiff in the premises so purchased could be constitutionally and legally impaired and rendered null and void.''
\end{quote}
\noindent --- Chief Justice Marshall, Opinion of the Court, 10 U.S. at 139.

\begin{quote}
``Is the clause to be considered as inhibiting the state from impairing the obligation of contracts between two individuals, but as excluding from that inhibition contracts made with itself? The words themselves contain no such distinction. They are general, and are applicable to contracts of every description.''
\end{quote}
\noindent --- Chief Justice Marshall, Opinion of the Court, 10 U.S. at 137.

\end{document}
